\section{Finite local identification of translation classes}
\label{sec:fingerprints}
An oriented bridge is an ordered pair of physical bodies, their shortest
segment of length $g\le R_+$, and one choice of common transverse
unit tangent. Clearance is not needed in this section. Translate its
source contact to zero and rotate the directed bridge to the positive
horizontal axis. The two facing boundaries are
\[
 (-\psi_0(y),y),\qquad (g+\psi_1(y),y),\qquad
 \psi_b(0)=\psi_b'(0)=0.
\]
Choose $r$ smaller than their common analytic radius. Define
\begin{equation}\label{eq:signature}
 F_K=\left(g,\left\{\frac{r^j}{j!}\psi_b^{(j)}(0):
                         b=0,1,\ 2\le j\le K\right\}\right).
\end{equation}
The norm is the maximum norm after fixing units from the prior bounds.
Simultaneous transverse reversal acts by multiplying the degree-$j$
coordinates by $(-1)^j$; denote this involution by $\iota$.
The ordering of source and opposite body is not quotiented out.

\begin{lemma}[A fixed-presentation compactness chart]
\label{lem:presentation-compactness}
The family of tables and bounded oriented bridges can be covered by a
finite disjoint union of compact parameter spaces on which every fixed
contact derivative is continuous and translation-class equivalence is
constant on each pair of discrete bridge-label strata. No reduced
lattice basis or half-open fundamental cell is required.
\end{lemma}
\begin{proof}
For each $1\le r\le r_0$, keep all allowed labelled presentations
$L$ with the prescribed bounds on $L$ and $L^{-1}$. Write the Steiner
centers as $Lt_i$, with $t_i\in[0,1]^2$, and describe each centered
body by its support function. Use Euclidean convergence for $L,t_i$
and uniform convergence on every closed narrower complex strip for the
support functions. Uniform holomorphic bounds give compactness in this
topology, hence convergence of every fixed derivative on the real circle.
Intersect with the closed physical and quantitative-margin conditions.
The resulting finite union over $r$ is compact. This is a redundant
parameter space, not a discontinuous choice of a reduced representative.

Normalize the source translate of an ordered bridge to $C_i$. Its target
is $C_j+Lk$. Since the two centers are within $b$ before translation,
$|k|\le\|L^{-1}\|(b+R_++2D_0)$. There are finitely many labels
$(i,j,k,\varepsilon)$, $\varepsilon\in\{+,-\}$, excluding the
identical-body case. For each label, the shortest pair of points is
unique: two distinct minimizing pairs would contradict strict convexity
by taking their midpoints. Compactness of the bodies then proves
continuity of these minimizing points. The positive contact Hessian
in \eqref{eq:normal-hessian} and the implicit-function theorem give
continuity of the normal graphs through every fixed derivative on a
smaller common disk. The transverse frame is chosen as $\varepsilon$
times the positive quarter-turn of the directed normal, so it has no
angular branch cut.

In any one such presentation, two labels are translation-equivalent
exactly when all four entries agree. Shape separation identifies the
types; invertibility of $L$ distinguishes their translates; uniqueness
of contacts and the sign distinguish the frames. These facts remain
true at limits because the positive margins remain. Replacing $C_i$
by $C_i+Lm_i$ merely relabels $k$ by $k+m_i-m_j$; replacing $L$
by $LU$ relabels $k$ by $U^{-1}k$. They preserve equivalence. At
fundamental-cell faces or lattice-reduction walls we keep both
presentations instead of selecting one discontinuously. This proves the
asserted chart formulation and justifies subsequences with fixed labels.
\end{proof}

\begin{lemma}[Separation of coexisting oriented bridges]
\label{lem:finite-separation}
There are $K_*<\infty$ and $\chi>0$, depending only on the compact
class and $R_+$, such that two oriented bridges of any one table have
$\|F_{K_*}-F'_{K_*}\|_\infty\ge\chi$, unless one is a translate
of the other with the same transverse orientation. In particular
$\|F_{K_*}-\iota F_{K_*}\|_\infty\ge\chi$ for every bridge.
\end{lemma}
\begin{proof}
First consider equality at every order. Analyticity gives equality of
both contact graphs on an interval. Two connected compact embedded
analytic curves which share an arc have the same image: the set of
points of local agreement is open, and accumulation of agreeing arcs
in a local graph gives agreement by the analytic identity theorem, so
it is closed as well. This applies to both boundaries in their channel
frames. Thus equality of all coefficients and of the gap produces one
Euclidean isometry between the complete ordered pair images.

By \eqref{eq:shape}, the two source bodies have the same type. Two
copies of this type differ by a translation, and their matching is
unique because the type has trivial Euclidean isometry group. The pair
isometry is therefore that translation. It maps the target body and the
contacts too. Its transverse orientations agree. In particular, reversing
one orientation cannot preserve all coefficients: that would give a
reflection symmetry of the source body.

We justify uniform finite order rather than merely pointwise separation,
using Lemma~\ref{lem:presentation-compactness}.
Take bounded lattice presentations and choose representative centers in
closed fundamental parallelograms. An ordered bridge is then specified
by $(i,j,k)$, where its source has type $i$ and its target is $C_j+Lk$.
The gap, body diameters, presentation and inverse-presentation bounds
bound $|k|$ uniformly. There are only finitely many such labels and
two transverse choices. The number $r$ takes finitely many values.
Unique shortest contacts of separated strictly convex bodies vary
continuously with the bodies, by compactness and uniqueness; their
normal-coordinate charts vary with every fixed derivative on a smaller
analytic disk.

Suppose the conclusion failed. For each integer $m$ choose one admissible
table and two inequivalent oriented bridges whose descriptors through
order $m$ differ by less than $1/m$. Pass to a subsequence with fixed
$r$, fixed discrete bridge labels and transverse choices, and convergent
bodies and lattice presentation. Compact analytic bounds give convergence
of every fixed contact coefficient. The limit table retains the separation,
shape and symmetry margins. The limit bridges agree at every order, so
by the preceding argument their discrete labels and orientations are
equal. This contradicts the fixed inequivalent labels. Distinct target
lattice copies cannot merge in the limit, by the positive separation and
inverse-lattice bounds. This proves that some finite order has a positive
uniform separation. Reducing it gives the stated $\chi$.
\end{proof}

This lemma is not finite-jet determination of an infinite-dimensional
analytic family. Its comparison is only between finitely many bounded
bridge orbits \emph{coexisting in one physical table}. A numerical
implementation needs a certified usable pair $(K_*,\chi)$ for its
prior class. The compactness proof supplies existence, not a general
bit-complexity bound for computing sharp prior constants. Conservative
class constants may be used. No actual fingerprint or complete shape is
part of those constants. Section~\ref{sec:effective} strengthens this
existence argument by an explicit numerical calibration and a value-only
fingerprint; no independently supplied $(K_*,\chi)$ is needed there.

\begin{proposition}[A robust registry and coherent local signs]
\label{prop:registry}
Assume every measured descriptor differs from its exact value by at most
$\varepsilon_f<\chi/16$. A registry stores the first measured descriptor
of each observed class. Compare a new descriptor to each stored one under
the two actions $1,\iota$, with threshold $\chi/2$.
Exactly one class and one sign match when that class has already appeared;
otherwise a new class is inserted. This conclusion is uniform over any
number of observations and over arbitrary, possibly nonperiodic, errors
within the certified bounds.
\end{proposition}
\begin{proof}
Correctly oriented measurements of the same class differ by at most
$2\varepsilon_f<\chi/8$. Measurements of distinct oriented classes,
including the opposite sign of one class, differ by at least
$\chi-2\varepsilon_f>7\chi/8$. Thus the comparison bands are disjoint.
Induction proves that the stored list has exactly one representative of
each observed unoriented class and that every returned relative sign is
correct. No transitive chaining of ambiguous similarities is used.
\end{proof}

The same rule evaluates a recorded class at an arbitrary newly encountered
copy. In exact local coordinates it is translation invariant: translating
contacts and flights changes neither the descriptor nor intrinsic offsets.
It therefore constructs a periodic gate without locating a fundamental
cell. At finite coordinate accuracy the class and sign decisions remain
exact, but bin assignments can change near cell boundaries. We will charge
that error explicitly rather than calling the measured coordinate map
itself exactly periodic.

The analytic and asymmetric hypotheses have distinct roles. Analyticity
extends agreement of contact graphs to complete images; shape separation
identifies the source type; asymmetry removes a nontranslational matching
of two copies of that type. Smooth jets alone do not justify the first
step. Dropping these hypotheses requires a different gate-identification
argument; the earlier registered theorems for larger classes are not
replaced by this lemma.
